\chapter*{Prefacio}
\addcontentsline{toc}{chapter}{Prefacio}

Este libro es la presentación pedagógica y autocontenida del sistema
\textbf{sotFS}: un filesystem cuya estructura interna no es un árbol de
inodos, como en \texttt{ext4} o \texttt{XFS}, sino un grafo dirigido
tipado, donde cada operación POSIX se modela como una reescritura algebraica
(DPO, \emph{double pushout}) y donde las propiedades de seguridad,
consistencia y crash-safety se enuncian como invariantes formales y se
verifican parcialmente con TLA\textsuperscript{+} y Coq.

\section*{Para quién es este libro}

Tres lectores en mente, en orden de comodidad creciente con el material:

\begin{itemize}
  \item Un \emph{ingeniero de sistemas} con base en filesystems POSIX,
    storage e idea de qué hace \texttt{fsync}, pero sin background previo
    en teoría de categorías ni verificación formal.
  \item Un \emph{estudiante de posgrado} en sistemas o métodos formales,
    interesado en cómo se baja un modelo matemático a una implementación
    Rust real, con tests, fuzzers, benchmarks y CI.
  \item Un \emph{practitioner de seguridad} interesado en provenance,
    monitoreo estructural (curvatura de Ricci-Ollivier sobre el grafo
    del FS) y proyecciones engañosas como mecanismo defensivo.
\end{itemize}

\section*{Cómo está organizado}

Catorce capítulos organizados en tres bloques:

\begin{description}
  \item[\sffamily Núcleo (caps. 1–5).] Motivación, preliminares
    matemáticos, el TypeGraph, hashing de contenido y reescritura DPO de
    todas las operaciones POSIX.
  \item[\sffamily Implementación y persistencia (caps. 6–9).] Cómo se baja
    el modelo a Rust seguro y concurrente; persistencia con \redb;
    capabilities y delegación; transacciones GTXN y \emph{crash refinement}.
  \item[\sffamily Verificación, monitoreo y exposición (caps. 10–14).]
    Specs TLA\textsuperscript{+}, pruebas en Coq con su deuda explícita,
    métricas estructurales (treewidth y curvatura), el \emph{functor de
    decepción} \(D_{\mathrm{FS}}\) y el adaptador \fuse{} que expone todo
    al kernel.
\end{description}

Cinco apéndices (catálogo de invariantes, tablas DPO, resultados TLC,
setup reproducible, soluciones a ejercicios ★ y ★★) cierran el volumen.

\section*{Convenciones}

\begin{itemize}
  \item \textbf{Cajas tipadas} sirven para destacar material verificable o
    peligroso. Las \emph{invariantes} (en azul) son propiedades que el
    sistema mantiene siempre. Las \emph{trampas} (en rojo) son errores
    históricos del proyecto, con su \emph{lección} explícita; muchas
    están enlazadas a un commit. Los \emph{ejercicios} (en ámbar)
    aparecen al final de cada capítulo.

  \item \textbf{Algoritmos} se presentan en pseudo-código estilo CLRS con
    \texttt{algorithm2e}, numerados por capítulo. Su correspondencia
    con código Rust real se cita por archivo y línea, p.~ej.\
    \citaLine{sotfs-graph/src/graph.rs}{571}.

  \item \textbf{Citas al repositorio} usan macros: \citaTLA{spec.tla},
    \citaCoq{archivo.v}, \citaCrate{nombre} y \citaCommit{sha}. Las
    citas académicas siguen el estilo numérico estándar.

  \item \textbf{Ejercicios graduados.} \dificultad{$\star$} memorística;
    \dificultad{$\star\star$} aplicación; \dificultad{$\star\star\star$}
    diseño o extensión. El apéndice~\ref{ap:soluciones} cubre los
    primeros dos niveles; los \dificultad{$\star\star\star$} llevan
    sólo una pista.
\end{itemize}

\section*{Lo que este libro no hace}

\begin{itemize}
  \item No reemplaza la lectura del código. Documenta y motiva la
    arquitectura de la versión \texttt{v0.2.5}; la fuente de verdad
    sigue siendo el repositorio.
  \item No oculta las debilidades formales. Hay cinco lemas
    \texttt{Admitted} en Coq (\(4\) en \texttt{DpoRmdir.v}, \(1\) en
    \texttt{DpoUnlink.v}); el cap.~\ref{cap:coq} los lista uno a uno.
  \item No cubre el crate \texttt{sotfs-experimental}: es código
    aspiracional, no integrado al runtime estable.
\end{itemize}

\section*{Agradecimientos}

A las herramientas que hicieron este libro posible: \LaTeX{} y
\texttt{algorithm2e}, \texttt{TikZ}, \texttt{biblatex}, TLA Tools,
\texttt{coqc} y \texttt{coqdoc}, \texttt{cargo-fuzz} y \texttt{cargo-llvm-cov}.
A la comunidad de teoría de categorías aplicada por insistir en que un
\emph{pushout} es la cosa más útil que uno aprende en la vida.

\bigskip
\noindent\hfill \emph{Buenos Aires, \today}
\clearpage
